\section{Conclusion} \label{sec:conclusion}
In this project, we applied statistical tools from econophysics to analyse the motion of bubbles in a coarsening two-dimensional foam, using Brownian motion as a baseline of pure randomness to determine the particular type of randomness of the bubble dynamics.\\

The bubble paths were compared with a Brownian walk of the same amount of time steps. While the Brownian walk looks erratic and uncorrelated, the bubble path looks to have more structured motion. Small clustered steps when the foam is relaxed, and large correlated jumps when the foam is restructuring. This inspired further investigation into the underlying statistics.\\

The mean squared displacement for Brownian walks follows the linear relation $\langle r^2(\tau) \rangle = 2D\tau$. The MSD for the bubbles were analysed in hopes to find the same relation. While the bubble MSD seems to follow an upward trend, the more bubbles disappear, the less data we have to average over, so it deviates a lot. It was hard to draw conclusions from this so we moved on to other statistical methods.\\

The displacement distribution deviate significantly from the Gaussian distribution of the Brownian motion, a sharp peak around the centre that tapers off. However, when we look at larger time lags $\Delta t$, the displacement distribution turns into a Gaussian distribution as the displacements become essentially random. More interestingly, for larger $\Delta t$ in the log returns distribution the central peak seems to persist, hinting at an underlying correlation in the log returns.\\

The CCDF analysis showed that the bubble distributions do not scale according to a Lévy distribution, unlike the scaling of the financial log returns. Instead of being linear, it was found that the function with two free parameters, $f(C, \Delta t) = c_1(C) \Delta t / (1 + c_2(C)\Delta t)$, describes the scaling of the CCDF cuts pretty well. A best fit was used on $c_1$ vs C and $c_2$ vs C, to get a closed form; $c_1 = 0.83(C+1)$, $c_2 = 0.0031\,C^{-0.61}$. Using these to obtain a fully collapsed master plot remains an goal for future work.\\

The analysis of the autocorrelation functions showed some interesting results. While the bubble displacements showed no significant correlation, as expected, the log returns ACF drops close to zero quickly however remains positive for the entire simulation. The volatility ACF, following the behaviour of the Dow Jones Industrial Average volatility ACF, follows a power decay law of $\tau^{-0.4}$, suggesting the bubble volatility share some statistical structure with financial markets.\\

Future work could explore the scaling of the CCDF further. A closed relation of the free parameters gives hope that they could be used in some way to collaps the cut scaling curve and use a feature of this collapsed plot such as the slope to collapse the original distributions onto a masterplot. Extending the analysis of the foams to different foam samples of different lengths, numbers of bubbles, and liquid fractions would also help to establish how general our findings are.